%%%%%%%%%%%%%%%%%%%%%%%%%%%%%%%%%%%%%%%%%%%%%%%%%%%%%%%%%%%%%%%%%%%%%%
% writeLaTeX Example: Academic Paper Template
%
% Source: http://www.writelatex.com
% 
% Feel free to distribute this example, but please keep the referral
% to writelatex.com
% 
%%%%%%%%%%%%%%%%%%%%%%%%%%%%%%%%%%%%%%%%%%%%%%%%%%%%%%%%%%%%%%%%%%%%%%
% How to use writeLaTeX: 
%
% You edit the source code here on the left, and the preview on the
% right shows you the result within a few seconds.
%
% Bookmark this page and share the URL with your co-authors. They can
% edit at the same time!
%
% You can upload figures, bibliographies, custom classes and
% styles using the files menu.
%
% If you're new to LaTeX, the wikibook is a great place to start:
% http://en.wikibooks.org/wiki/LaTeX
%
%%%%%%%%%%%%%%%%%%%%%%%%%%%%%%%%%%%%%%%%%%%%%%%%%%%%%%%%%%%%%%%%%%%%%%
\documentclass[twocolumn,showpacs,%
  nofootinbib,aps,superscriptaddress,%
  eqsecnum,prd,notitlepage,showkeys,10pt]{revtex4-1}

\usepackage{amssymb}
\usepackage{amsmath}
\usepackage{graphicx}
\usepackage{dcolumn}
\usepackage{hyperref}

\begin{document}

\title{Your Paper}
\author{Jinwon Cho, Adam Yonge, and Carrie Farberow}
\affiliation{National Laboratory of the Rockies}

\begin{abstract}
bafasdfasdfasdf.
\end{abstract}

\maketitle

\section{Introduction}
The transition to sustainable energy infrastructure requires efficient methods to produce and deliver clean fuels, making hydrogen a vital resource for decarbonization and achieving a net-zero emission society.1, 2 Given the logistical challenges associated with storing and transporting gaseous hydrogen in large volumes, chemical carriers such as ammonia (NH3), methane, or methanol are frequently used.1-4 Among these, NH3 is an effective carbon-free hydrogen carrier, although its on-demand conversion to hydrogen requires a significant energy input.5, 6
Autothermal reforming (ATR) offers an advantageous pathway for NH3 decomposition, as it facilitates a series of internal chemical processes that generate the requisite energy for the overall reaction to proceed.7-9 However, the ATR process presents a mechanistic challenge in achieving a delicate balance between generating heat, primarily through ammonia combustion, and maximizing hydrogen selectivity.10-12 Experimental observations using supported ruthenium oxide (RuO2) nanoparticles in packed-bed reactors reveal highly dynamic operating conditions, including a significant temperature gradient from 900°C at the entrance to 300°C at the outlet, following the initiation of the process.13, 14 This extensive thermal and chemical environment, where oxygen is completely consumed at the outlet, leads to numerous surface compositions where ammonia reactions are possible.
Consequently, the selective outcome of ATR, specifically the favorability of desired products (H2 and N2) versus undesired oxidized products (H2O and NO/NO2), is hypothesized to be fundamentally dictated by the surface coverage of oxygen on the RuO2 and Ru catalyst.15-17 Mechanistically, product selectivity is controlled by the interaction of intermediates with the surface oxygen. The formation of desired products, such as N2, leads to the retention of an active site (∗), whereas the formation of undesired oxidized species consumes surface oxygen atoms, resulting in the creation of oxygen vacancies (∗v).15
While previous computational and experimental studies have provided insights into ammonia reactions over RuO2 surfaces, these investigations often concentrate on the most commonly observed RuO2 (110) cleavage, focusing primarily on surface coverages characterized by equal amounts of Ru and O sites. To address the challenge of understanding selectivity across the broad range of compositions observed in ATR, our study focuses on the RuO2(110) surface, alongside of Ru(0001) surface, across the full spectrum of relevant oxygen compositions, including O-Rich, O/Ru equivalent, and Ru-Rich states.
Specifically, we performed the first-principles calculations of RuO2 (110) and Ru (0001) surfaces to explore the essential link between surface oxygen availability and selective ATR. Our analysis focuses on the competitive reaction steps: 1) NH3 adsorption and the initial N−H cleavage; 2) the competitive formation of H2O versus H2 desorption from adsorbed atomic hydrogen (H∗) and 3) the competition between NO desorption and the formation of the desired product, N2, from adsorbed atomic nitrogen (N∗).

\section{Calculation Details}
The calculations were performed using spin-polarized DFT within the Perdew−Burke−Ernzerhof (PBE) functional, as implemented in the Vienna Ab Initio Simulation Package (VASP). The projector augmented wave (PAW) method with a planewave basis set was employed to describe the interaction between the core and valence electrons. An energy cutoff of 430 eV was applied for the planewave expansion of the electronic eigenfunctions. A (3 × 3 × 1) Monkhorst–Pack mesh of k points was used to calculate the geometries and total energies by Brillouin zone integration [42]. A supercell slab was built for the Ir/Ta2O5 catalyst, which is composed of orthorhombic δ-Ta2O5 (3 × 3) surface unit cells with four atomic layers, and a face-centered cubic (FCC) Ir(111) (2 × 2) bilayer slab was constructed, as displayed in Fig. 7a. The slab was separated from its periodic images in the vertical direction by a vacuum space corresponding to seven atomic layers. The bottom two layers were fixed at the corresponding bulk positions, whereas the upper four layers were fully relaxed using the conjugate gradient method until the residual forces on all the constituent atoms became smaller than 5 × 10−2 eV Å−1. The reaction free energy  for each elementary OER step was defined as the difference between the free energies of the initial and final states:(1)where , , , and  denote the difference in reaction energy, the change in zero-point energy (ZPE), the change in entropy, and the change in free energy due to the applied electrode potential, respectively. ZPE corrections were calculated from the vibrational frequencies according to standard methods at 298 K, and those of gas-phase molecules were obtained from thermodynamic databases [43].  is defined as(2)where , , and  are the number of protons–electrons transferred in pairs, the elementary charge, and the applied electrode potential versus standard hydrogen potential, respectively.
\label{sec:examples}

\subsection{Sections}

Use \texttt{section}s and \texttt{subsection}s to organize your document. \LaTeX{} handles all the formatting and numbering automatically. Use \texttt{ref} and \texttt{label} for cross-references --- this is Section~\ref{sec:examples}, for example.

\subsection{Tables and Figures}

Use \texttt{tabular} for basic tables --- see Table~\ref{tab:widgets}, for example. You can upload a figure (JPEG, PNG or PDF) using the files menu. To include it in your document, use the \texttt{includegraphics} command (see the comment below in the source code).

% Commands to include a figure:
%\begin{figure}
%\includegraphics[width=\textwidth]{your-figure's-file-name}
%\caption{\label{fig:your-figure}Caption goes here.}
%\end{figure}

\begin{table}
\centering
\begin{tabular}{l|r}
Item & Quantity \\\hline
Widgets & 42 \\
Gadgets & 13
\end{tabular}
\caption{\label{tab:widgets}An example table.}
\end{table}

\subsection{Mathematics}

\LaTeX{} is great at typesetting mathematics. Let $X_1, X_2, \ldots, X_n$ be a sequence of independent and identically distributed random variables with $\text{E}[X_i] = \mu$ and $\text{Var}[X_i] = \sigma^2 < \infty$, and let
$$S_n = \frac{X_1 + X_2 + \cdots + X_n}{n}
      = \frac{1}{n}\sum_{i}^{n} X_i$$
denote their mean. Then as $n$ approaches infinity, the random variables $\sqrt{n}(S_n - \mu)$ converge in distribution to a normal $\mathcal{N}(0, \sigma^2)$.

\subsection{Lists}

You can make lists with automatic numbering \dots

\begin{enumerate}
\item Like this,
\item and like this.
\end{enumerate}
\dots or bullet points \dots
\begin{itemize}
\item Like this,
\item and like this.
\end{itemize}

\begin{acknowledgments}

We thank\dots

\end{acknowledgments}

\end{document}